\section{Related work}\label{sec:related}

\paragraph{Concurrent work.} Lu et al.~\citep{lu2026cachetrap} study the same phenomenon on SGLang with prefill--decode
disaggregation and a host KV tier, serving a hybrid mixture-of-experts model (Qwen3.6-35B-A3B) on Terminal-Bench and SWE-bench Pro
trajectories. On hardware, in single runs, as their run accounting implies, a ten-minute surge cut goodput by about 88\% over the following hour,
terminating the burst sessions at the end of the burst restored service, and in a run that let them finish the base population kept
growing after the last burst session left. Their fluid model, $\dot N=\lambda-X(N)/K$ with an empirical hit curve $h(W/C)$ of the
working set over the cache, gives the stable and unstable equilibria, the vulnerable window between the low-reuse departure rate
$D_{\mathrm{low}}$ and the peak, the statement that arrivals below $D_{\mathrm{low}}$ let the population drain, a resource condition
for cache-first vulnerability, and a drain-time estimate. In simulation they map vulnerability, show that changing post-surge
arrivals alone reverses recovery, and report cache hysteresis (a lower hit rate after collapse at a comparable working set). Their
controller, Basin\footnote{Their Basin controller is unrelated to our BASIN-1 and BASIN-2 tests, which probe the basin of
attraction of the good state.}, combines fixed session priority, admission by predicted peak context, decode-memory lending and
shedding; on H200 hardware an arrival-time prototype of it achieves 1.20--1.70 times Continuum's goodput, and the predictor-driven
design is evaluated in simulation. Against this, we add an exactly analysed idealization of the trapped state with its tests
(\S\ref{sec:trap}--\ref{sec:validation}), an engine test of the reading that the throughput measured in a trap exceeds the drain
threshold (Corollary~\ref{cor:trapthr}, proved in model~F, our idealization of an FCFS engine on which every turn misses; CUT-A and
CUT-B are consistent with it; their systems keep some reuse), a first-principles entry model with an
out-of-sample record (\S\ref{sec:entry}), the registered and replicated rescue of an established zero-hit trap at the same offered
load without shedding, with release and persistence (\S\ref{sec:control}), and registered matched-arm evidence on colocated engines,
including a second engine (\S\ref{sec:evidence}). They have real traces, disaggregated serving, a larger model, SLO goodput and a full
mitigation system, which we do not.

\paragraph{Thrashing and metastability.} The thrashing collapse and the working-set remedy are Denning's~\citep{denning1968thrashing,%
denning1968working}; metastability from a unimodal throughput curve and bistability under finite sources were analysed for
random-access channels and multiprogrammed systems~\citep{kleinrock1975aloha,carleial1975bistable,courtois1975decomposability,%
nelson1987stochastic}. Metastable failures were named by Bronson et al.~\citep{bronson2021metastable}, studied in production
incidents and modelled by Huang et al.~\citep{huang2022metastable}, whose trigger-and-control design our paired test follows, and
analysed with queueing models and Markov chains~\citep{isaacs2025analyzing,habibi2024msf,alvaro2025formal}; Dong poses metastability
in queues as an open problem~\citep{dong2022metastability}. These works characterize the bad regime as a state or an absorbing
region; as far as we found, none gives its growth law or the composition of its population. Service systems whose service slows in
congestion show related bistability~\citep{chan2014speedup,dong2015slowdowns}.

\paragraph{Overloaded and re-entrant queues.} Overloaded processor sharing has a Crump--Mode--Jagers branching
structure~\citep{grishechkin1992processor}, linear growth~\citep{jeanmarie1994transient} and a fluid limit whose shape converges to an
invariant one~\citep{puha2006overloaded}; overloaded multiclass FIFO queues with abandonment and no re-entry have fluid
limits~\citep{jennings2013fluid}, and overloaded networks have bottleneck fluid approximations~\citep{chen1991discrete}. Re-entrant
lines~\citep{kumar1993reentrant} and their fluid models~\citep{dai1996stability,bramson1996convergence} are studied mostly for
stability, and for convergence to equilibria at load at most one. Closest to model~F, Khanin et al.\ study steady, constant-flow solutions of fluid FIFO re-entrant lines,
in which an overloaded station's queue grows linearly~\citep{khanin2001steady}; at a single station their fixed point is our Perron
root, and a closing remark covers fractional exits, which is our random turn count. FIFO queues with Bernoulli or deterministic
feedback are classical in the stable regime~\citep{takacs1963feedback,vandenberg1989deterministic}. Round robin with feedback
to the tail of the queue, which model~F is at the level of sessions, was analysed in equilibrium for time-shared
systems~\citep{kleinrock1967timeshared,coffman1968feedback}, and overloaded polling systems have a random fluid limit with a branching
structure~\citep{remerova2014random}. In heavy traffic, multiclass FIFO feedback queues show state-space collapse~\citep{reiman1988multiclass,dai1995multiclass}, of
which Theorem~\ref{thm:trap}(b) is the fluid counterpart. Theorem~\ref{thm:trap} adds the renewal recursion along FIFO orbits, a
discrete renewal equation of Lotka--Leslie type~\citep{feller1968introduction}, and, with it, convergence from an arbitrary initial
state with an exact virtual origin. Theorem~\ref{thm:premium} is the visit-indexed analogue of the overloaded processor-sharing growth
law, whose departures can exceed the capacity~\citep{jeanmarie1994transient}; its turn-count and cost-growth factors have no
counterpart we found. Corollary~\ref{cor:trapthr} is the operational reading of Theorem~\ref{thm:premium} with the classical drain
condition, and its processor-sharing counterpart follows from~\citep{jeanmarie1994transient} in the same way.

\paragraph{Queueing models of LLM serving.} Ao et al.\ give fluid-guided scheduling with memory constraints~\citep{ao2025wait} and show
that, in a decode-only model without prefix reuse, the eviction-free fixed point can be unstable~\citep{ao2026service}; Dai et al.\ prove
throughput optimality of work-conserving policies without cache state~\citep{dai2025throughput}; Nie et al.\ give stability conditions
under KV constraints for one-shot requests~\citep{nie2026queueing}; Dexter et al.\ analyse FCFS and longest-prefix-match scheduling under
prefix reuse~\citep{dexter2025prefix}; and Yao et al.\ derive a mean-field LRU hit ratio for multi-turn sessions with exogenous gaps
between turns~\citep{yao2026multiturn}, which our entry model extends with hit-dependent running reservations
(\S\ref{sec:entry}).

\paragraph{KV management for agents.} Agent-serving systems keep, offload or evict the KV that sessions leave idle in tool calls,
order requests to preserve reuse, and bound the number of active agents: PEEK~\citep{xie2026peek}, Continuum~\citep{li2025continuum},
Sutradhara~\citep{biswas2026sutradhara}, CONCUR~\citep{chen2026concur}, ThunderAgent~\citep{kang2026thunderagent} and its Dynamo
port~\citep{dhanani2026dynamo}, EfficientAgent~\citep{shao2026efficientagent}, AgentKV~\citep{anon2026agentkv},
MISA-T~\citep{hong2026misat}, MARS~\citep{wang2026mars}, VAMP~\citep{kim2026vamp} and KAIROS~\citep{yuan2026kairos}, among others
(Appendix~\ref{app:supp}). Several report hit-rate collapse as concurrency grows~\citep{chen2026concur,yuan2026agentic,hong2026misat,%
kim2026vamp}, describe the loop between eviction, cold prefill and queueing for closed or bounded populations, or assert that a short
overload can leave an instance persistently inefficient~\citep{yuan2026kairos}. We claim none of these mechanisms; a session cap is
classical load control~\citep{denning1976optimal,cherkasova2002session}.
